\begin{figure*}[t]
  \centering
  \begin{tikzpicture}[
    font=\fontsize{8}{9}\selectfont,
    node distance=5mm and 7mm,
    box/.style={draw=black!55, rounded corners=1.5pt, minimum height=7mm,
                minimum width=22mm, align=center, line width=.6pt},
    stage/.style={box, fill=blue!7},
    evidence/.style={box, fill=orange!10},
    flow/.style={-{Latex[length=2mm]}, line width=.8pt, draw=black!75}
  ]
    \node[box] (input) {Input\\$x$};
    \node[stage, right=of input] (encode) {1. Encode\\$h=f_\theta(x)$};
    \node[stage, right=of encode] (reason) {2. Mechanism\\$z=g_\phi(h)$};
    \node[box, right=of reason] (output) {Output\\$\hat y$};
    \node[evidence, below=of reason] (test) {Matched test\\prediction $\rightarrow$ result};
    \draw[flow] (input) -- node[above, font=\fontsize{7}{8}\selectfont]{$x$} (encode);
    \draw[flow] (encode) -- node[above, font=\fontsize{7}{8}\selectfont]{$h$} (reason);
    \draw[flow] (reason) -- node[above, font=\fontsize{7}{8}\selectfont]{$z$} (output);
    \draw[flow] (reason) -- (test);
  \end{tikzpicture}
  \caption{\textbf{Method overview and evidence interface.} The diagram names the
  input, each reusable component, its mathematical object, the output, and the
  experiment that tests the proposed mechanism.}
  \label{fig:method}
\end{figure*}
